\subsection{Text equivalent and evidence key for the Ark plates}

The first plate presents a rough-pencil, three-quarter study of a rectangular
wooden chest.  It marks the biblical proportions of two and one-half cubits
long by one and one-half cubits wide and high without converting the cubit to
a modern measurement.  The chest is acacia---the Douay and Vulgate tradition's
setim wood---overlaid with gold inside and outside.  A separate pure-gold cover,
the mercy seat or propitiatory, rests above it.  Two stylized outlines mark the
hammered-gold cherubim at the cover's ends, their wings overshadowing it and
their faces turned toward it.  Four rings receive two wood poles overlaid with
gold; the poles are shown remaining in the rings.  The cherubim's bodies and
the chest's ornament are explicitly not reconstructed (Ex 25:10--22;
37:1--9).

The diagram does not flatten the sources concerning the contents.  Exodus and
Deuteronomy place the covenant tablets inside.  The Pentateuch relates the
manna vessel and Aaron's rod to the Testimony and puts the law-book beside the
Ark.  Hebrews 9:4 retrospectively gathers the manna vessel, rod, and tablets
within its sanctuary description.  At the Temple dedication, 1 Kings 8:9 says
that only the two tablets were inside.  No transfer or disappearance of the
other items is invented.

\begin{fourstable}{Setting}{What Scripture attests}{What the plate shows}{What remains unknown}
Wilderness Tabernacle & The commanded Tent is raised; the Ark and cover are
placed behind the veil; cloud covers the Tent and glory fills it (Ex 26;
40:18--21, 34--38 in modern numbering). & A schematic court, Tent, Holy Place,
veil, Most Holy Place, and Ark. & The plate is not to scale and claims no
surviving structure, exact cubit conversion, or archaeological footprint.\\
Shiloh & The Tent of Meeting is established at Shiloh; later language names a
house or temple of the Lord, doors, a lamp, and the Ark (Jos 18:1; 1 Sm 1:9;
3:3, 15; 4:3--4). & A dashed sanctuary boundary and an Ark marker, with the
uncertainty printed inside the panel. & The terms do not recover a floor plan,
building material, sequence of architectural change, or the Ark's place within
a known footprint.\\
David's tent on Zion & David pitches a tent, brings the Ark into it, and
appoints worship before it; the older Tabernacle and altar remain at Gibeon
(2 Sm 6:17; 1 Chr 16:1, 39--40). & A deliberately symbolic tent outline with
the Ark within. & The form and location of David's tent are unknown; the symbol
does not identify it with the Mosaic Tent or a later Temple building.\\
Solomon's Temple & The Ark is brought up from the City of David into the inner
sanctuary, beneath the great cherubim; only the tablets are inside; cloud and
glory fill the house (1 Kgs 6:19--28; 8:1--11; 2 Chr 5:2--14). & A simplified
porch, Holy Place, inner sanctuary, Ark, and overarching cherubim. & This is a
textual teaching plan, not a claim to recover the precise First-Temple
footprint, furnishings in elevation, or the lost Ark.\\
\end{fourstable}

These settings trace a real history without turning every sanctuary into the
same structure.  They also discipline the Marian fulfillment.  Mary is not a
fifth sacred building: she is the free, living mother in whom the Son truly
takes flesh.  The old veil, chamber, cloud, and Ark become intelligible in
their Christological fullness when the Spirit overshadows her and the divine
Word dwells personally within her.
